\begin{table}[h!]
    \centering
    \renewcommand{\arraystretch}{1.5} % Tăng khoảng cách dòng cho thoáng
    \caption{Bảng phân công vai trò và nhiệm vụ chi tiết của nhóm}
    \label{tab:team_assignment}
    
    % Cấu trúc cột: 
    % p{2.5cm}: Cột Thành viên cố định rộng 2.5cm
    % p{4cm}: Cột Vai trò cố định rộng 4cm
    % X: Cột Nhiệm vụ tự động giãn hết phần còn lại của trang
    \begin{tabularx}{\textwidth}{@{} p{2.5cm} p{4cm} X @{}}
    \toprule
    \textbf{Thành viên} & \textbf{Vai trò chuyên môn} & \textbf{Nhiệm vụ chi tiết \& Kết quả thực hiện} \\ 
    \midrule
    
    \textbf{Huỳnh Quốc Huy} 
    & Xử lý dữ liệu \& Tích hợp AI Tạo sinh 
    & 
    \textbf{1. Tiền xử lý dữ liệu (Data Preprocessing):}
    \begin{itemize}[leftmargin=*, nosep, after=\vspace{0.3em}]
        \item Thu thập, làm sạch và chuẩn hóa các nguồn dữ liệu đa phương thức (hình ảnh, văn bản, thông số kỹ thuật).
        \item Xây dựng pipeline xử lý dữ liệu tự động để phục vụ quá trình huấn luyện.
    \end{itemize}
    \textbf{2. Xây dựng Hệ thống Demo (Web Application):}
    \begin{itemize}[leftmargin=*, nosep]
        \item Thiết kế và phát triển giao diện người dùng (Frontend) cho Web Demo.
        \item Xây dựng Backend API để kết nối giữa giao diện và các mô hình gợi ý, hiển thị kết quả trực quan cho người dùng.
    \end{itemize}
    \textbf{3. Tích hợp LLM (Explainability):}
    \begin{itemize}[leftmargin=*, nosep]
        \item Nghiên cứu và tích hợp mô hình ngôn ngữ lớn để sinh văn bản giải thích lý do gợi ý.
        \item Thiết kế kỹ thuật \textit{Prompt Engineering} để tạo ra các giải thích tự nhiên và phù hợp ngữ cảnh.
    \end{itemize}
    \\ \midrule
    
    \textbf{Nguyễn Gia Huy} 
    & Mô hình hóa Sản phẩm \& Nghiên cứu Nâng cao 
    & 
    \textbf{Triển khai mô hình Cross-Encoder:}
    \begin{itemize}[leftmargin=*, nosep]
        \item Xây dựng và tinh chỉnh (fine-tune) kiến trúc Cross-Encoder chuyên biệt cho giai đoạn xếp hạng lại (Re-ranking).
        \item Đánh giá hiệu năng và phân tích mức độ cải thiện độ chính xác của hệ thống sau khi tích hợp Cross-Encoder so với các phương pháp truy xuất thô.
    \end{itemize}
    \\ \midrule
    
    \textbf{Đàm Đức Huy} 
    & Mô hình hóa Người dùng
    & 
    \textbf{1. Xây dựng Tháp Sản phẩm (Item Tower):}
    \begin{itemize}[leftmargin=*, nosep, after=\vspace{0.3em}]
        \item Thiết kế và cài đặt các bộ mã hóa đặc trưng (\texttt{ImageEncoder}, \texttt{TextEncoder}, \texttt{TabularEncoder}).
        \item Hiện thực hóa cơ chế \textit{Item Fusion} để tạo vector đại diện sản phẩm thống nhất.
    \end{itemize}
    \textbf{2. Xây dựng Tháp Người dùng (User Tower):}
    \begin{itemize}[leftmargin=*, nosep, after=\vspace{0.3em}]
        \item Cài đặt kiến trúc \textit{Transformer/SASRec} để mô hình hóa chuỗi hành vi tuần tự của người dùng.
        \item Hiện thực hóa các kỹ thuật \textit{Masked Mean Pooling} và \textit{Positional Encoding}.
    \end{itemize}
    \\ \bottomrule
    \end{tabularx}
\end{table}

% \newpage

% \begin{table}[H]
%     \centering
%     \caption{Kế hoạch thực hiện đề tài}
%     \label{tab:project-timeline}
%     \renewcommand{\arraystretch}{1.3}
%     \begin{tabularx}{1.0\linewidth}{lX}
%     \toprule
%     \textbf{Tuần (Thời gian)} & \textbf{Nội dung thực hiện} \\
%     \midrule
%     w37 (08/09 -- 14/09) & Tìm hiểu tổng quan, xác định phạm vi đề tài và dataset. \\
%     w38 (15/09 -- 21/09) & Thu thập, phân tích và làm sạch bộ dữ liệu \textit{Amazon Reviews}. \\
%     w39 (22/09 -- 28/09) & Thực hiện embedding dữ liệu sản phẩm và người dùng. \\
%     w40 (29/09 -- 05/10) & Xây dựng và cài đặt Item Tower (Tích hợp Fusion Gate). \\
%     w41 (06/10 -- 12/10) & Xây dựng User Tower với kiến trúc SASRec để mô hình hóa chuỗi hành vi. \\
%     w42 (13/10 -- 19/10) & Huấn luyện mô hình Two-Tower (Loss function: \textit{InfoNCE} + In-batch Negatives). \\
%     w43 (20/10 -- 26/10) & Triển khai cơ chế Retrieval và chiến lược Hybrid Reranking. \\
%     w44 (27/10 -- 02/11) & Tối ưu hóa mô hình: Tinh chỉnh Hyperparameters ($\alpha, \beta$, Learning rate) để cải thiện chỉ số NDCG/Recall. \\
%     w45 (03/11 -- 09/11) & Xây dựng Demo (FastAPI/ReactJS) và tích hợp nhanh API LLM để giải thích kết quả. \\
%     w46 (10/11 -- 16/11) & Đánh giá toàn diện hệ thống: So sánh với các baseline và phân tích trường hợp lỗi (Error Analysis). \\
%     w47 (17/11 -- 23/11) & Tổng hợp kết quả thực nghiệm, vẽ biểu đồ so sánh hiệu năng. \\
%     w48 (24/11 -- 30/11) & Viết báo cáo chương 1--3 (Lý thuyết, Thiết kế hệ thống). \\
%     w49 (01/12 -- 07/12) & Viết báo cáo chương 4--5 (Thực nghiệm, Kết luận, Hướng phát triển tương lai). \\
%     w50 (08/12 -- 14/12) & Rà soát lần cuối và hoàn thiện slide thuyết trình. \\
%     \bottomrule
%     \end{tabularx}
% \end{table}